\documentclass[conference]{IEEEtran}
\usepackage[T1]{fontenc}
\usepackage[utf8]{inputenc}
\usepackage{amsmath,amssymb}
\usepackage{array}
\usepackage{booktabs}
\usepackage{multirow}
\usepackage{url}
\usepackage{microtype}
\newcolumntype{P}[1]{>{\raggedright\arraybackslash}p{#1}}
\setlength{\emergencystretch}{2em}

\title{Where Do We Put Excess Carbon Dioxide?\\A Permanence-Aware Portfolio for Atmospheric Removal and Storage}
\author{Research Plan Author Agent\\Xcientist four-agent simulation}

\begin{document}
\maketitle

\begin{abstract}
The question of where to put excess carbon dioxide (CO2) is not solved by identifying a single large sink. A scientifically useful formulation asks which portfolio can deliver additional, measurable, durable, and equitable net atmospheric CO2 removal while accounting for lifecycle emissions, energy, land, water, materials, ecological effects, leakage, reversal, monitoring, and long-term stewardship. This report proposes PERSIST-CO2, a Permanence, Energy, Reversal, Storage Integrity, and Trade-offs framework. The Survey stage finds that geological, biological, mineral, direct-air-capture, bioenergy, enhanced-weathering, and ocean-based pathways have different retention timescales and constraints. The Idea stage selects a mixed portfolio over geology-first, biology-first, engineered-removal-first, and nominal-cost routes. The ExperimentDesign stage specifies pathway-specific stock dynamics, durability-adjusted mass balance, resource and infrastructure constraints, monitoring capacity, robust optimization, controls, ablations, and falsification. The direct design claim is that geological or mineral storage should anchor durable carbon management, while bounded biological and engineered pathways provide complementary benefits only when their reversal and lifecycle liabilities are priced. This is a design-only proposal: it reports no measured removal, storage integrity, or deployment result.
\end{abstract}

\begin{IEEEkeywords}
carbon dioxide removal, carbon sequestration, geological storage, biological carbon, permanence, leakage, monitoring and verification, robust optimization
\end{IEEEkeywords}

\section{Introduction}
The amount of CO2 emitted by the energy and industrial system is much larger than can be managed by placing a few projects at a few sites. The central engineering question is therefore not simply where a molecule can be injected or planted. It is whether a pathway produces a verified net atmospheric benefit after the energy, materials, transport, monitoring, leakage, reversal, and stewardship required to keep that benefit. This distinction is essential because captured mass, removed mass, and retained atmospheric benefit are different quantities.

The broad question is reframed as follows: which portfolio of carbon dioxide removal and storage pathways can deliver additional, measurable, durable, and equitable net atmospheric CO2 removal at meaningful scale while minimizing lifecycle energy, land, water, material, ecological, leakage, and reversal risks? The proposed answer is PERSIST-CO2 (Permanence, Energy, Reversal, Storage Integrity, and Trade-offs for CO2).

The IPCC Working Group III assessment identifies carbon dioxide removal (CDR) as a necessary element of net-zero pathways for counterbalancing residual emissions from hard-to-transition sectors, while emphasizing that methods differ in process, storage timescale, maturity, potential, cost, co-benefits, adverse effects, and governance requirements~\cite{ipccwg3}. The National Academies research agenda similarly frames CDR and reliable sequestration as a research problem involving benefits, risks, sustainable scale, commercial viability, and specific technical questions across terrestrial, coastal, engineered, and geological methods~\cite{nas2019}. These sources support a portfolio experiment, not an assumption that any nominal capacity is immediately deployable.

The report makes three contributions. First, it defines a common accounting boundary that separates gross capture, lifecycle burden, retained stock, leakage, reversal, and net atmospheric benefit. Second, it proposes an optimization experiment in which geological and mineral durability, biological reversibility, engineered energy demand, monitoring capacity, and regional burdens are explicit. Third, it defines a decision rule for answering the practical question: use durable geological or mineral storage as the anchor for concentrated and atmospheric removal, use biological and coastal pools where their co-benefits justify bounded reversible storage, and reject a portfolio that reaches a tonne target only by hiding lifecycle or stewardship liabilities.

The work follows the four-agent workflow. Survey maps evidence and gaps. Idea selects PERSIST-CO2 from competing routes. ExperimentDesign specifies equations, scenarios, controls, and falsification. Author freezes those outputs into this IEEE-style proposal. There are no observed results and no claim that any pathway has already removed a stated amount of CO2.

\section{Evidence-Constrained Survey}
\subsection{Why storage locations are not interchangeable}
Geological storage places compressed CO2 in deep saline formations, depleted reservoirs, or reactive formations. Its durability potential is high relative to many surface pools, but the benefit depends on site characterization, injectivity, pressure management, caprock and well integrity, plume migration, induced seismicity, monitoring, and post-closure stewardship. Mineralization can make carbon part of a stable mineral phase, but reaction kinetics, rock supply, crushing, transport, water, and energy constrain scale.

Biological storage includes forests, soils, harvested wood, wetlands, mangroves, and seagrasses. These pools can be built with biodiversity, erosion, water, or livelihood co-benefits, but their carbon stocks can be reversed by wildfire, drought, heat, pests, harvest, land-use change, or policy failure. A biological tonne is not equivalent to a geological tonne unless the accounting horizon, disturbance distribution, and replacement obligation are stated.

Engineered atmospheric removal includes direct air capture with storage (DACCS), bioenergy with carbon capture and storage (BECCS), enhanced weathering, and ocean-based approaches. They may address diffuse atmospheric CO2 or provide durable storage, but their net value depends on low-carbon energy, land, water, minerals, transport, ecological effects, and verification. The IPCC reports that CDR methods differ substantially in readiness, potential, costs, side effects, and governance requirements~\cite{ipccwg3}.

\subsection{The scale and research bottleneck}
The U.S. Energy Information Administration maintains international CO2 emissions data, which provides a useful emissions-scale context but does not itself establish a removal target or a storage guarantee~\cite{eia}. The key inference is qualitative: a global emissions inventory must be connected to a portfolio model that includes both gross emissions avoidance and residual removal. A historical emissions number cannot be subtracted from a nominal storage capacity without a common time, boundary, and verification protocol.

The National Academies report identifies unanswered questions needed to assess benefits, risks, sustainable scale, commercial viability, and the research and development required for CDR and reliable sequestration~\cite{nas2019}. Its scope includes land and coastal ecosystem management, accelerated weathering, BECCS, DAC, geological sequestration, and related approaches. This supports a design in which MRV, scale, and stewardship are modeled as capacity constraints rather than added after a cost curve has been optimized.

The IPCC Synthesis Report places mitigation, adaptation, and residual risk in a coupled climate system~\cite{ipccsyr}. For carbon storage, the implication is that an apparently successful global tonne count can still create local harm if land is diverted from food, water use rises in a stressed basin, mineral extraction damages ecosystems, or liability is shifted to communities after project finance ends.

\begin{table*}[t]
\centering
\caption{Evidence anchors and claims permitted in the PERSIST-CO2 design.}
\label{tab:evidence}
\small
\begin{tabular}{P{2.8cm}P{4.3cm}P{4.7cm}P{3.2cm}}
\toprule
Evidence anchor & Directly supported proposition & Design use & Forbidden inference \\
\midrule
IPCC WGIII Ch. 12~\cite{ipccwg3} & CDR counterbalances residual emissions in net-zero pathways; methods differ in durability, maturity, potential, cost, effects, and governance & Pathway-specific retention, lifecycle, resource, and governance parameters & CDR permits continued unrestricted fossil emissions \\
National Academies~\cite{nas2019} & Benefits, risks, sustainable scale, commercial viability, and R\&D questions remain important & MRV, scale, and stewardship are modeled as research variables and constraints & Nominal capacity is deployable capacity \\
EIA international data~\cite{eia} & International CO2 emissions data provide scale context & Declare inventory year, boundary, gases, and uncertainty before optimization & A historical inventory is a verified current removal target \\
IPCC AR6 Synthesis~\cite{ipccsyr} & Mitigation and residual risk interact across systems and regions & Report distributional burden, ecological co-impacts, and residual risk & A global tonne target guarantees local climate benefit \\
\bottomrule
\end{tabular}
\end{table*}

\subsection{Gap ledger}
The Survey Agent accepts six connected gaps. First, pathways are not consistently compared using durability-adjusted lifecycle accounting. Second, leakage and reversal timescales are often reduced to a single discount factor. Third, monitoring, reporting, and verification are treated as a reporting layer instead of a limited operational capacity. Fourth, energy, land, water, mineral, transport, and ecological limits interact nonlinearly at scale. Fifth, global targets can hide land tenure, food security, water, biodiversity, and environmental-justice burdens. Sixth, stewardship and liability can outlive project finance, but closure and remediation triggers are under-specified.

The survey rejects four overclaims. It rejects the claim that every captured tonne is permanently stored. It rejects the equivalence of forests, soils, coastal pools, and geological reservoirs without a declared horizon and reversal policy. It rejects future CDR as permission to delay gross emissions reductions. It rejects nominal pore volume, biomass area, or modeled technical potential as proof of socially accepted, monitored, and deployable capacity.

\section{Idea Synthesis}
\subsection{Competing directions}
A geology-first route sends most eligible CO2 to saline formations or depleted reservoirs. It has a strong durability hypothesis, but can be brittle when site characterization, injection pressure, pipeline capacity, monitoring, or public acceptance is limiting. A biology-first route emphasizes forests, soils, wood, wetlands, mangroves, and seagrasses. It can deploy quickly and offer co-benefits, but is vulnerable to reversal, additionality disputes, land competition, and climate feedbacks.

An engineered-removal-first route scales DACCS, BECCS, mineralization, enhanced weathering, or ocean approaches by nominal cost or modeled potential. It can target atmospheric CO2 and may offer durable storage, but it can shift emissions upstream through energy demand and create material, water, land, ecological, or governance burdens. An avoidance-only route prevents storage liabilities and should be a control, but cannot address residual emissions from hard-to-transition activities or legacy atmospheric CO2.

\subsection{Primary direction: PERSIST-CO2}
PERSIST-CO2 combines the strengths of these routes while refusing to average away their differences. Its four layers are:
\begin{enumerate}
\item carbon accounting for gross capture, lifecycle burden, leakage, reversal, and retained benefit;
\item reservoir dynamics for pathway-specific retention and failure processes;
\item resource and infrastructure planning for energy, land, water, minerals, transport, wells, MRV, and stewardship;
\item portfolio governance for additionality, staged deployment, liability reserves, monitoring triggers, and credit retirement.
\end{enumerate}

The central hypothesis is:
\begin{quote}
A permanence-aware mixed portfolio anchored by durable geological or mineral storage will have lower worst-case climate-benefit regret than a single-pathway or nominal-cost strategy under plausible energy, land, water, leakage, reversal, monitoring, and governance futures.
\end{quote}

The hypothesis is ambitious but falsifiable. It does not claim that geology always wins or that biological carbon is unimportant. It predicts that a portfolio should allocate each pathway to the role for which its physical and governance properties are favorable, and should not pay a permanent credit for a reversible stock without a replacement or liability mechanism.

\subsection{Predictions}
PERSIST-CO2 predicts that durability-adjusted accounting will change pathway rankings relative to nominal cost. It predicts that MRV capacity can become binding before nominal pore volume, biomass area, or theoretical atmospheric removal potential is exhausted. It predicts that the optimal robust portfolio will combine early gross emissions avoidance with durable storage and bounded biological or engineered removal, rather than maximize captured tonnes through one route. It predicts that staged deployment with monitoring triggers will outperform delayed or unbounded scale-up when reversal or energy uncertainty is material.

\section{ExperimentDesign}
\subsection{Study object and scenario ensemble}
The study object is a staged carbon-management portfolio, not a forecast of one technology. A portfolio contains gross emissions reductions, pathway shares, siting decisions, transport and injection capacity, MRV intensity, stewardship reserves, and trigger rules. Regions are represented by industrial sources, energy systems, land and water accounts, candidate storage sites, ecological zones, and exposed access groups. The planning horizon contains milestone decisions and finer operational periods for capture, transport, injection, restoration, and monitoring.

The scenario ensemble varies the emissions baseline, energy carbon intensity and availability, capture and removal costs, retention and leakage, fire and drought, pests and land-use change, geological integrity and pressure, mineral supply, land and water scarcity, transport and injection build rates, MRV coverage and error, monitoring latency, governance capacity, and social acceptance. Scenarios are split into development, stress-test, and held-out sets before optimization. A scenario is a stress test, not a prediction.

\subsection{Durability-adjusted carbon accounting}
Let $M_{i,t}$ denote gross CO2 captured or removed by pathway $i$ at time $t$, $p_i(t)$ the fraction retained at the declared evaluation time, $l_i(t)$ the fraction lost through leakage or reversal, and $e_i$ the lifecycle CO2-equivalent burden per unit gross removal. The net atmospheric benefit is:
\begin{equation}
R_t^{\mathrm{net}}=\sum_i M_{i,t}\left[p_i(t)-l_i(t)-e_i\right].
\label{eq:netremoval}
\end{equation}
All terms are converted to a common CO2-equivalent boundary. The experiment reports gross capture, lifecycle emissions, retained stock, leakage, reversal, and net benefit separately. When $e_i$, $p_i$, or $l_i$ is uncertain, the uncertainty is propagated rather than replaced by a favorable point estimate.

For a storage pool $S_i(t)$, stock dynamics are:
\begin{equation}
\frac{dS_i(t)}{dt}=I_i(t)-\lambda_i(t)S_i(t)-V_i(t),
\label{eq:stock}
\end{equation}
where $I_i(t)$ is injected or accumulated carbon, $\lambda_i(t)$ is the leakage rate, and $V_i(t)$ is disturbance or reversal flux. Geological pathways use pressure, plume, well, and caprock failure modes. Biological pathways use disturbance, harvest, and land-use transition modes. Ocean and weathering pathways use dissolution, transport, alkalinity, and ecological uncertainty modes. The equation is a bookkeeping interface, not a claim that all pathways have the same physics.

\subsection{Resource and infrastructure constraints}
For each region $r$ and time $t$, energy consumed by capture, compression, transport, removal, and MRV is bounded by low-carbon energy after priority demand:
\begin{equation}
\sum_i a_{i,r,t}M_{i,r,t}+E^{\mathrm{MRV}}_{r,t}
\leq E^{\mathrm{clean}}_{r,t}-E^{\mathrm{priority}}_{r,t},
\label{eq:energy}
\end{equation}
where $a_i$ is pathway energy intensity, $E^{\mathrm{MRV}}$ is monitoring energy, and $E^{\mathrm{priority}}$ protects existing essential demand. Analogous constraints cover eligible land, water withdrawal, mineral throughput, pipeline capacity, wells, monitoring crews, data latency, and stewardship reserves.

Additionality is tested against a declared baseline land or industrial trajectory. A project receives no removal credit for carbon that would have been stored without the intervention. Transport emissions, parasitic energy, leakage outside the project boundary, and displaced land-use emissions remain in the common accounting boundary. Biological deployment is bounded by food, biodiversity, tenure, and water constraints; geological deployment is bounded by site-specific injectivity, pressure, plume, and monitoring limits.

\subsection{Robust portfolio optimization}
Let $z$ be a staged portfolio and $J_w(z)$ a documented scalarization of cost, negative net benefit, energy burden, land and water burden, ecological impact, leakage risk, MRV shortfall, and distributional burden in scenario $w$. The minimax-regret policy is:
\begin{equation}
\begin{aligned}
z^*=\mathop{\mathrm{arg\,min}}\limits_{z\in\mathcal{Z}}\;&\max_{w\in\mathcal{W}}\left|J_w(z)-J_w^*\right|\nonumber\\
\text{subject to }&R_t^{\mathrm{net}}\geq \underline{R}_t,\nonumber\\
&\operatorname{MRVShortfall}(z,w)\leq \overline{m},\nonumber\\
&\operatorname{Burden}_{r}(z,w)\leq \overline{b}_r,
\end{aligned}
\label{eq:robust}
\end{equation}
where $\mathcal{Z}$ is the feasible portfolio set, $\mathcal{W}$ is the scenario set, $J_w^*$ is the best feasible scenario-specific value, $\underline{R}_t$ is the minimum net benefit, and $\overline{m}$ and $\overline{b}_r$ are monitoring and regional burden limits. The design reports Pareto fronts and sensitivity to weights, horizons, accounting boundaries, and retention assumptions.

\subsection{Monitoring and stewardship}
Every project registers a baseline, mass-flow measurement, retention horizon, uncertainty interval, leakage pathway, responsible operator, financial reserve, closure plan, and credit-retirement rule. Geological monitoring includes pressure, seismic, geochemical, plume, and well-integrity indicators as appropriate. Biological monitoring includes biomass, soil, disturbance, land-use, fire, drought, harvest, and additionality indicators. Engineered pathways include energy, materials, process emissions, transport, and product or mineral fate.

A trigger is activated when plume migration, pressure, ecosystem loss, fire, drought, reversal, or measurement uncertainty crosses a predeclared envelope. Responses can pause injection, retire credits, fund remediation, increase monitoring, or move future volume to another pathway. Hysteresis prevents noisy repeated switching. A trigger identifies a mandatory review; it cannot bypass legal authority or assign liability without an accountable institution.

\begin{table}[t]
\centering
\caption{Pathway roles and binding risks to be tested.}
\label{tab:pathways}
\footnotesize
\begin{tabular}{P{1.55cm}P{2.1cm}P{2.5cm}}
\toprule
Pathway & Intended role & Binding risks \\
\midrule
Geological & Durable anchor for concentrated CO2 and DACCS & Site integrity, pressure, wells, transport, MRV, stewardship \\
Mineral & Durable solid storage & Rock supply, energy, water, kinetics, mining impacts \\
Biological & Co-benefits and bounded near-term removal & Fire, drought, land, additionality, reversal \\
Coastal & Ecosystem-linked storage & Disturbance, area, permanence, coastal change, MRV \\
DACCS & Atmospheric removal where clean energy exists & Energy, cost, materials, water, siting \\
BECCS & Energy plus removal in selected systems & Land, food, biomass, supply chain, lifecycle emissions \\
Weathering and\\ ocean & Exploratory scale options & Verification, side effects, transport, governance \\
\bottomrule
\end{tabular}
\end{table}

\subsection{Controls and ablations}
Controls include avoidance-only, geology-first, biology-first, engineered-removal-first, and a nominal-cost optimizer that ignores durability. A perfect-MRV control estimates an information upper bound but is not treated as deployable. Ablations remove durability accounting, MRV capacity, energy and resource constraints, equity constraints, or stewardship triggers one at a time. The incremental value of each layer is measured by changes in net benefit, regret, resource use, reversal exposure, monitoring shortfall, and regional burden.

\subsection{Analysis sequence}
The analysis freezes the emissions boundary, evaluation horizon, additionality baseline, pathway parameter ranges, scenario split, weights, thresholds, and reliability requirements. It then constructs the ensemble, calibrates pathway and infrastructure models, optimizes all controls and PERSIST-CO2, evaluates held-out scenarios and extreme events, runs ablations, and performs sensitivity and model-discrepancy analyses. The final report presents retained carbon by pathway, net atmospheric benefit, leakage and reversal, lifecycle burden, resource use, MRV shortfall, stewardship reserve, Pareto fronts, and distributional outcomes.

\section{Expected Interpretations and Falsification}
The experiment has four pre-specified interpretation classes. A \textit{durable-portfolio outcome} occurs if PERSIST-CO2 reduces held-out worst-case regret and meets net-benefit, MRV, resource, and burden constraints. A \textit{durability-reversal outcome} occurs if biological or nominal-cost rankings deteriorate sharply when disturbance and retention horizons are represented. An \textit{MRV-bound outcome} occurs if verification capacity, not physical storage, limits the deployable portfolio. A \textit{single-pathway outcome} occurs if one pathway matches or exceeds the mixed portfolio across held-out scenarios without hidden resource or stewardship exclusions.

\begin{table}[t]
\centering
\caption{Pre-registered interpretation rules.}
\label{tab:decisions}
\footnotesize
\begin{tabular}{P{1.6cm}P{2.25cm}P{2.25cm}}
\toprule
Outcome & Scenario signature & Interpretation \\
\midrule
Durable portfolio & Regret falls and constraints hold & Mixed, permanence-aware design supported \\
Reversal & Nominal biological advantage disappears & Reversal accounting is decision-critical \\
MRV bound & Monitoring capacity binds first & Verification infrastructure is a scale limit \\
Single pathway & One route wins under equal constraints & Portfolio diversity is not justified in tested domain \\
\bottomrule
\end{tabular}
\end{table}

PERSIST-CO2 is falsified if it does not outperform single-pathway and nominal-cost controls on held-out durability-adjusted regret under the same accounting and physical constraints. It is also falsified if its advantage disappears under held-out reversal or energy scenarios, requires perfect monitoring or unbounded land and water, or meets a tonne target only by excluding stewardship and remediation costs. The framework does not retain a pathway as a decision component if its additional volume never changes net benefit, risk, or a monitored action.

A negative result is useful. If geology is limited by monitoring or site integrity, the model identifies that bottleneck. If biology performs well under a declared horizon and disturbance regime, the result supports a bounded biological role rather than permanent equivalence. If DACCS fails because of energy, the result specifies the energy condition that must change. No negative result is converted into a claim that all CDR is ineffective.

\section{Decision Policy and Sensitivity}
The portfolio decision is evaluated at more than one temporal horizon. A 20-year horizon tests whether a project delivers near-term atmospheric benefit and avoids rapid reversal. A century-scale horizon tests whether the same credit remains valid after disturbance, leakage, replacement, and stewardship events. Geological retention is not assumed to be perfect at either horizon, and biological retention is not discarded merely because it is reversible; instead, each pathway receives a crediting rule that matches its measured retention and replacement obligation.

The model reports sensitivity to five choices that can otherwise dominate the result. The first is the lifecycle boundary, including upstream energy, transport, displaced production, and construction. The second is the retention horizon and the shape of the pathway-specific survival curve. The third is the treatment of disturbance events and correlated failures, such as drought and wildfire. The fourth is the availability of low-carbon energy and the competition between removal and essential demand. The fifth is the social and ecological burden weight assigned to land, water, minerals, and access groups. A portfolio is robust only when its qualitative role survives these changes, not when one favorable parameter set produces a large number.

Monitoring has option value because it can prevent invalid credits and redirect future deployment. The experiment therefore compares a fixed monitoring budget with an adaptive budget that increases at a trigger. The value of extra MRV is measured by the reduction in expected credit retirement, remediation cost, and worst-case regret. A monitoring investment is retained only if it changes a decision, detects a failure earlier than the stewardship lead time, or narrows uncertainty enough to alter the portfolio. Data volume without a decision effect is not counted as climate benefit.

The resulting policy rule is deliberately operational. First, reduce gross emissions and reserve durable storage for residual emissions and verified atmospheric removal. Second, route concentrated streams to characterized geological or mineral sites when energy, transport, pressure, and monitoring constraints are satisfied. Third, use biological and coastal pools for bounded removal and co-benefits when additionality, tenure, disturbance, and replacement obligations are explicit. Fourth, permit DACCS, BECCS, weathering, or ocean pathways to expand only when their lifecycle and ecological conditions are met. Finally, retire credits or substitute pathways when measured retention leaves the declared envelope. This rule is a model output to be tested, not a legal allocation of storage rights.

\section{Data Contract and Evaluation Metrics}
The design requires a data contract before any optimizer is allowed to rank pathways. The contract makes every tonne traceable to a source, a process boundary, a time interval, a location, a measurement method, and a responsible institution. A pathway record contains the gross-flow variable, energy and material inputs, transport distance, storage-pool identifier, retention model, leakage and reversal modes, uncertainty distribution, MRV method, land and water accounts, ecological indicators, and closure obligation. Missing values are not silently set to zero. They are labeled as missing, bounded with an interval, or excluded from the decision until a justified prior is recorded.

The same contract is applied to avoided emissions and removal. Avoidance is measured against an explicit counterfactual production or energy trajectory, whereas removal requires evidence that atmospheric CO2 entered the managed carbon stock. A captured stream that would otherwise have been emitted is not automatically atmospheric removal, and a biomass stock that would have existed without the project is not additional removal. The optimizer therefore receives separate ledgers for avoided emissions, gross captured carbon, gross atmospheric removal, lifecycle emissions, retained carbon, leakage, reversal, and credit retirement. This prevents a favorable label from substituting for a mass balance.

Evaluation uses metrics that retain physical meaning. The primary metric is held-out worst-case regret in durability-adjusted net atmospheric benefit subject to the same constraints for every comparator. Secondary metrics are cumulative net benefit, retained fraction at 20 years and at the century-scale horizon, lifecycle emissions per net tonne, clean-energy use per net tonne, land and water occupation, mineral throughput, MRV shortfall, credit-retirement volume, remediation reserve, and regional burden. A strategy that delivers more gross capture but less retained benefit, or that shifts burden to an unrepresented region, cannot win by a single nominal-cost score.

The experiment also records the shape of failure. For each scenario, the first binding constraint is tagged as energy, site integrity, transport, MRV, land, water, mineral supply, ecological burden, equity, or stewardship. The time from an indicator crossing its trigger to effective remediation is compared with the estimated time for additional leakage or reversal. This lead-time metric distinguishes a monitoring system that merely documents damage from one that can still preserve the credited climate benefit. Correlated failures are retained: a drought can reduce biomass, raise electricity demand, increase wildfire probability, and lower hydropower availability in the same scenario.

\begin{table}[t]
\centering
\caption{Minimum record required before a pathway receives a modeled credit.}
\label{tab:datacontract}
\footnotesize
\begin{tabular}{P{1.85cm}P{2.35cm}P{2.0cm}}
\toprule
Record & Required fields & Decision consequence if absent \\
\midrule
Flow & Gross mass, time, location, instrument, uncertainty & No net-tonne credit \\
Boundary & Energy, materials, transport, displaced effects & Use interval or exclude \\
Storage & Pool, retention curve, leakage modes, horizon & Credit capped at evidence \\
Additionality & Counterfactual and land-use baseline & Avoided emissions only \\
MRV & Indicator, frequency, latency, coverage, error & Increase reserve or pause \\
Responsibility & Operator, closure plan, finance, trigger & No durable classification \\
Co-impacts & Land, water, ecology, access groups & Apply regional constraint \\
\bottomrule
\end{tabular}
\end{table}

For uncertainty analysis, the design uses nested sampling. Outer draws represent structural alternatives, such as a retention curve family, lifecycle boundary, or ecosystem disturbance model. Inner draws represent parameter uncertainty conditional on that structure. Each policy is evaluated on identical draws so that differences are attributable to the decision rule rather than to unequal randomness. The result is reported as a distribution and a scenario envelope, with a clear distinction between parameter uncertainty, model discrepancy, and unobserved future conditions.

The computational comparison is staged to expose where the conclusion comes from. First, all pathways are evaluated using a common unit and no optimization, producing a transparent accounting dashboard. Second, each control is optimized under a shared feasibility set. Third, PERSIST-CO2 is optimized with durability, monitoring, and stewardship layers activated. Fourth, the same policies are replayed after the held-out scenario labels are hidden. Finally, stress events are injected after the policy is selected. This last replay tests whether a policy can respond to failure without retrospectively changing its crediting rule.

The minimum evidence for accepting the primary hypothesis is a consistent reduction in held-out worst-case regret across the declared sensitivity families, accompanied by a nonnegative improvement in net atmospheric benefit and no violation of energy, resource, regional burden, MRV, or stewardship constraints. If the mixed portfolio wins only under one horizon, one lifecycle boundary, or one weight vector, the result is reported as conditional rather than robust. If a simpler route wins under the same constraints, PERSIST-CO2 is rejected for that domain and its proposed hierarchy is revised. This evaluation standard makes a strong portfolio claim possible while keeping the mechanism of failure visible.

\section{Implementation, Governance, and Reproducibility}
The project starts with a versioned registry. Each pathway, project, site, measurement, emissions factor, retention curve, uncertainty interval, lifecycle boundary, and governance commitment receives an identifier and update history. The registry records open, restricted, delayed, missing, and quality-flagged observations. Results are not compared across versions without preserving the assumptions that generated them.

The computational stack has a data layer, pathway and reservoir models, an infrastructure and resource layer, a robust optimizer, and a decision layer. The decision layer reports scenario envelopes and trigger status rather than issuing an automatic credit or legal mandate. Synthetic scenarios and test cases are included so that the mass balance, stock dynamics, optimization, and ablations can be reproduced without restricted site data.

The deployment protocol separates three commitments. Immediate commitments are gross emissions reductions, measurement continuity, additionality baselines, and low-regret monitoring. Option-preserving commitments are modular capture, transport corridors, wells, restoration capacity, and interoperable MRV. Triggered commitments are expansion, crediting changes, remediation, or pathway substitution after a predeclared threshold. This structure prevents a future removal promise from becoming an excuse to defer current emissions reduction.

Environmental safeguards are hard constraints where data permit and otherwise explicit co-impact reports. A low-operational-carbon pathway is not assumed sustainable if it consumes scarce water, competes with food, displaces communities, damages biodiversity, increases mining impacts, or transfers closure liability. Actual deployment requires site characterization, permitting, local consultation, independent verification, and long-term institutional responsibility beyond the simulation.

\section{Risks and Scope}
The largest scientific risk is model discrepancy. Reservoir, ecosystem, atmospheric, energy, and lifecycle models may disagree. PERSIST-CO2 uses multiple structures and treats disagreement as decision-relevant uncertainty. A recommendation that survives only one simulator is not robust.

The largest accounting risk is leakage across system boundaries. Imported biomass, displaced land use, upstream energy, transport, product substitution, and uncertain removal permanence can make a project appear better than its system effect. Gross emissions, cumulative emissions, net benefit, removal dependence, leakage, and reversal are therefore reported separately.

The largest governance risk is the gap between a crediting period and the duration of storage responsibility. The registry records responsible operators, financial reserves, closure conditions, remediation triggers, and credit retirement. If those fields cannot be assigned, the pathway is not counted as verified durable removal even if its physical potential is large.

This report is a research proposal, not a global storage atlas, a site approval, or an empirical estimate of how many tonnes can be removed. It reports no observed storage integrity, no measured lifecycle result, and no deployed portfolio. The proposed scientific claim is narrower and actionable: the credible place for excess CO2 is a hierarchy of sinks in which gross emissions are reduced first, durable geological or mineral reservoirs anchor verified storage, and biological or engineered pathways are used only inside explicit energy, resource, reversal, monitoring, equity, and stewardship limits.

\section{Conclusion}
There is no single container large enough to answer the question “where do we put all the excess CO2?” The answer must be a portfolio with a common accounting boundary. Geological and mineral storage are the strongest candidates for durable anchoring, especially for concentrated sources and atmospheric removal, but they require site-specific characterization, infrastructure, monitoring, and post-closure responsibility. Forests, soils, wood products, wetlands, mangroves, and seagrasses can contribute co-benefits and bounded removal, but their reversibility must be priced. DACCS, BECCS, weathering, and ocean approaches can expand the option set when their energy, material, ecological, and governance conditions are satisfied.

PERSIST-CO2 makes this hierarchy testable. It computes net atmospheric benefit after lifecycle emissions, leakage, and reversal; tracks storage stocks over different horizons; treats MRV and stewardship as scarce capacities; constrains land, water, energy, minerals, and regional burdens; and selects policies using held-out worst-case regret. The strongest expected result is not a promise of unlimited removal. It is a decision rule: deploy gross emissions reductions immediately, build durable storage and verification where they are feasible, retain biological and engineered options as bounded complements, and retire or remediate credits when monitoring shows that permanence has failed.

The answer is therefore direct. We can manage excess atmospheric CO2 only by putting different fractions into different reservoirs under different contracts of permanence, measurement, and responsibility. PERSIST-CO2 should be accepted only if it demonstrates lower held-out regret and higher verified net benefit than simpler strategies without violating resource, ecological, equity, monitoring, or stewardship constraints. If it fails, the failure identifies which assumed sink cannot carry its proposed share.

\begin{thebibliography}{00}
\bibitem{eia} U.S. Energy Information Administration, ``International Carbon Dioxide Emissions Data,'' U.S. Department of Energy, 2026. [Online]. Available: \url{https://www.eia.gov/international/}.
\bibitem{ipccwg3} Intergovernmental Panel on Climate Change, ``Chapter 12: Cross sectoral perspectives,'' in \textit{Climate Change 2022: Mitigation of Climate Change}, Working Group III, 2022. [Online]. Available: \url{https://www.ipcc.ch/report/ar6/wg3/chapter/chapter-12/}.
\bibitem{nas2019} National Academies of Sciences, Engineering, and Medicine, \textit{Negative Emissions Technologies and Reliable Sequestration: A Research Agenda}. Washington, DC, USA: The National Academies Press, 2019. [Online]. Available: \url{https://www.nationalacademies.org/projects/DELS-BASCPR-16-01/publication/25259}.
\bibitem{ipccsyr} Intergovernmental Panel on Climate Change, \textit{Climate Change 2023: Synthesis Report}. Geneva, Switzerland: IPCC, 2023. [Online]. Available: \url{https://www.ipcc.ch/report/ar6/syr/}.
\end{thebibliography}

\end{document}
